\documentclass[11pt,a4paper]{article}

% ---------------- PACKAGES ----------------
\usepackage[utf8]{inputenc}
\usepackage[T1]{fontenc}
\usepackage[french]{babel}
\usepackage{lmodern}
\usepackage{geometry}
\usepackage{xcolor}
\usepackage{tcolorbox}
\usepackage{hyperref}
\usepackage{enumitem}
\usepackage{titlesec}
\usepackage{array}
\usepackage{longtable}
\usepackage{booktabs}
\usepackage{textcomp}
\usepackage{fancyhdr}

\tcbuselibrary{breakable}

% ---------------- PAGE ----------------
\geometry{margin=2.2cm}
\setlength{\parindent}{0pt}
\setlength{\parskip}{0.45em}
\setlist[itemize]{leftmargin=1.2cm, itemsep=0.2em}

% ---------------- COLORS ----------------
\definecolor{myblue}{RGB}{30,70,150}
\definecolor{mygreen}{RGB}{30,120,60}
\definecolor{myorange}{RGB}{200,120,40}
\definecolor{myred}{RGB}{170,50,50}
\definecolor{softgray}{RGB}{245,245,245}

% ---------------- HYPERREF ----------------
\hypersetup{
	colorlinks=true,
	linkcolor=myblue,
	urlcolor=myblue,
	citecolor=myblue
}

% ---------------- TITLES ----------------
\titleformat{\section}
{\Large\bfseries\color{myblue}}
{\thesection.}{0.5em}{}

\titleformat{\subsection}
{\large\bfseries\color{myorange}}
{\thesubsection.}{0.5em}{}

% ---------------- HEADER / FOOTER ----------------
\pagestyle{fancy}
\fancyhf{}
\lhead{Robot Auto-Équilibré STM32F4}
\rhead{Carnet de bord}
\cfoot{\thepage}
\lfoot{\small Pourvu qu'aujourd'hui soit mieux qu'hier.}

% ---------------- BOXES ----------------
\newtcolorbox{motivation}{
	breakable,
	colback=green!5,
	colframe=mygreen,
	title=Principe du projet,
	fonttitle=\bfseries
}

\newtcolorbox{sessionbox}{
	breakable,
	colback=blue!3,
	colframe=myblue,
	title=Journal de séance,
	fonttitle=\bfseries
}

\newtcolorbox{warningbox}{
	breakable,
	colback=red!4,
	colframe=myred,
	title=Point d'attention,
	fonttitle=\bfseries
}

\newtcolorbox{componentbox}{
	breakable,
	colback=orange!4,
	colframe=myorange,
	title=Fiche composant,
	fonttitle=\bfseries
}

% ---------------- DOCUMENT ----------------
\title{\textbf{Carnet de Bord\\Robot Auto-Équilibré STM32F4}}
\author{Ithiel Gabriel Djifa DENYIGBA}
\date{Version 1.0 -- Juillet 2026}

\begin{document}
	
	% ---------------- COVER ----------------
	\begin{titlepage}
		\centering
		\vspace*{2cm}
		
		{\Huge\bfseries\color{myblue} Carnet de Bord\\[0.3cm] Robot Auto-Équilibré STM32F4\par}
		
		\vspace{1.2cm}
		
		{\Large Développement d'un robot auto-équilibré capable de se déplacer et d'évoluer vers la navigation autonome.\par}
		
		\vspace{1.5cm}
		
		\begin{motivation}
			\centering
			\Large\textbf{Pourvu qu'aujourd'hui soit mieux qu'hier.}
			
			\vspace{0.3cm}
			Progresser chaque jour, même légèrement, est une victoire.
		\end{motivation}
		
		\vfill
		
		\begin{tabular}{>{\bfseries}l l}
			Auteur : & Ithiel Gabriel Djifa DENYIGBA \\
			Projet : & Robot auto-équilibré STM32F4 \\
			Version : & 1.0 \\
			Début : & Juillet 2026 \\
			Budget initial : & 200~\texteuro{} \\
		\end{tabular}
		
		\vspace{2cm}
	\end{titlepage}
	
	\tableofcontents
	\newpage
	
	% ================================
	% CHAPITRE 1
	% ================================
	\section{Pourquoi ce projet ?}
	
	Ce projet est un laboratoire personnel pour apprendre et mettre en pratique les notions fondamentales des systèmes embarqués, de l'automatique et de la robotique mobile.
	
	Il doit permettre de développer des compétences solides en :
	
	\begin{itemize}
		\item systèmes embarqués ;
		\item programmation STM32 ;
		\item électronique appliquée ;
		\item contrôle moteur ;
		\item automatique et PID ;
		\item traitement du signal ;
		\item fusion de capteurs ;
		\item robotique mobile ;
		\item architecture logicielle ;
		\item gestion de projet technique.
	\end{itemize}
	
	L'objectif n'est pas seulement de construire un robot. Le robot est le résultat visible. Le véritable objectif est de progresser comme ingénieur en apprenant à concevoir, tester, corriger, documenter et améliorer un système réel.
	
	\begin{motivation}
		Le robot n'est pas seulement une machine. Il devient un support d'apprentissage, un terrain d'expérimentation et une preuve concrète de progression.
		
		\centering
		\textbf{Pourvu qu'aujourd'hui soit mieux qu'hier.}
	\end{motivation}
	
	% ================================
	% CHAPITRE 2
	% ================================
	\section{Cahier des charges}
	
	\subsection*{Objectif général}
	
	Construire un robot auto-équilibré basé sur une carte STM32F4, capable de se stabiliser, de se déplacer et, à terme, d'évoluer vers une navigation autonome.
	
	\subsection*{Version 1 -- Stabilisation de base}
	
	Le robot doit :
	
	\begin{itemize}
		\item tenir debout en équilibre ;
		\item avancer ;
		\item reculer ;
		\item être contrôlé par une boucle temps réel ;
		\item exploiter une IMU pour estimer son inclinaison ;
		\item utiliser des moteurs avec encodeurs pour mesurer le mouvement.
	\end{itemize}
	
	\subsection*{Version 2 -- Robot connecté}
	
	Ajouter :
	
	\begin{itemize}
		\item Wi-Fi via ESP32 ;
		\item télémétrie vers PC ;
		\item réglage des paramètres PID à distance ;
		\item contrôle manuel depuis une interface distante.
	\end{itemize}
	
	\subsection*{Version 3 -- Navigation autonome}
	
	Ajouter :
	
	\begin{itemize}
		\item navigation d'un point A vers un point B ;
		\item estimation de déplacement par odométrie ;
		\item évitement d'obstacles ;
		\item planification de trajectoire simple.
	\end{itemize}
	
	\subsection*{Contraintes de conception}
	
	\begin{itemize}
		\item privilégier l'apprentissage progressif ;
		\item valider chaque brique séparément avant intégration ;
		\item garder une architecture logicielle claire ;
		\item documenter chaque séance ;
		\item éviter d'ajouter des fonctionnalités avancées avant stabilisation.
	\end{itemize}
	
	% ================================
	% CHAPITRE 3
	% ================================
	\section{Roadmap de développement}
	
	\begin{longtable}{p{3cm}p{5.5cm}p{5.5cm}}
		\toprule
		\textbf{Phase} & \textbf{Objectif} & \textbf{Compétences travaillées} \\
		\midrule
		Phase 1 & Prise en main STM32 & GPIO, UART, Timers, PWM, interruptions, ADC \\
		Phase 2 & Briques robot & IMU, I2C, moteur, encodeur, filtre complémentaire \\
		Phase 3 & Intégration mécanique & châssis, câblage, alimentation, centre de gravité \\
		Phase 4 & Stabilisation & PID angle, boucle temps réel, tuning \\
		Phase 5 & Déplacement & vitesse, odométrie, contrôle dynamique \\
		Phase 6 & Autonomie & Wi-Fi, télémétrie, navigation, évitement d'obstacles \\
		\bottomrule
	\end{longtable}
	
	\begin{motivation}
		Chaque phase est une marche. L'objectif n'est pas de sauter les étapes, mais de les valider proprement.
	\end{motivation}
	
	% ================================
	% CHAPITRE 4
	% ================================
	\section{Journal de bord}
	
	Chaque séance suit ce format :
	
	\begin{itemize}
		\item Objectif ;
		\item Réalisation ;
		\item Problèmes rencontrés ;
		\item Solution trouvée ;
		\item Apprentissage ;
		\item Prochaine étape.
	\end{itemize}
	
	\newpage
	
	%%%%%%%%%%%%%%%%%%%%%%%%%%%%%%%%%%%%%%%%%%%%%%%%%%%%%
	\subsection*{Séance 0 -- Naissance du projet}
	
	\begin{sessionbox}
		
		\textbf{Date :} Juillet 2026
		
		\vspace{0.3cm}
		
		\textbf{Objectif :}
		
		Définir la vision du projet, concevoir son architecture et sélectionner l'ensemble du matériel nécessaire à la réalisation du robot auto-équilibré.
		
		\vspace{0.3cm}
		
		\textbf{Réalisation :}
		
		\begin{itemize}
			\item Définition de la vision du projet.
			\item Création du carnet de bord.
			\item Élaboration de la feuille de route de développement.
			\item Définition des différentes phases du projet.
			\item Étude des architectures possibles.
			\item Comparaison des différentes cartes STM32.
			\item Sélection des moteurs, capteurs et composants.
			\item Vérification de la compatibilité de tous les éléments.
			\item Validation de l'architecture électrique.
			\item Commande de l'ensemble des composants.
		\end{itemize}
		
		\vspace{0.3cm}
		
		\textbf{Matériel commandé :}
		
		\begin{itemize}
			\item STM32 Nucleo-F446RE.
			\item Module IMU MPU6050.
			\item Driver moteur TB6612FNG.
			\item Deux motoréducteurs avec encodeurs.
			\item Batterie LiPo 2S 7,4 V -- 2200 mAh.
			\item Chargeur LiPo avec équilibrage.
			\item Convertisseur Buck LM2596.
			\item Connecteurs Deans avec câbles silicone.
			\item Interrupteur ON/OFF.
			\item Multimètre numérique.
			\item Cutter de précision.
			\item Plaques de PVC expansé.
			\item Fils de câblage.
			\item Câbles Dupont.
		\end{itemize}
		
		\vspace{0.3cm}
		
		\textbf{Budget du projet :}
		
		\begin{center}
			\Huge
			\textbf{200~\texteuro{}}
			
			\Large
			Version 1
		\end{center}
		
		\vspace{0.3cm}
		
		\textbf{Difficultés rencontrées :}
		
		\begin{itemize}
			\item Choisir entre plusieurs cartes STM32.
			\item Comprendre l'alimentation d'un robot mobile.
			\item Sélectionner une batterie compatible.
			\item Comprendre le rôle du Buck Converter.
			\item Choisir des moteurs adaptés au contrôle d'équilibre.
			\item Vérifier la compatibilité entre tous les composants.
		\end{itemize}
		
		\vspace{0.3cm}
		
		\textbf{Ce que j'ai appris :}
		
		Avant de programmer un robot, il faut d'abord le concevoir. Un bon projet commence par une architecture solide. Le choix des composants est aussi important que les algorithmes qui les feront fonctionner.
		
		\vspace{0.3cm}
		
		\textbf{Prochaine séance :}
		
		Installer STM32CubeIDE et commencer la prise en main de la STM32 avec le premier programme GPIO.
		
	\end{sessionbox}
	
	\begin{motivation}
		Aujourd'hui, je n'ai encore écrit aucune ligne de code.
		
		Pourtant, le projet est déjà né.
		
		Chaque grande réalisation commence par une décision.
		
		Aujourd'hui, j'ai décidé de construire quelque chose qui me fera progresser.
		
		\vspace{0.4cm}
		\centering
		\Large\textbf{Pourvu qu'aujourd'hui soit mieux qu'hier.}
	\end{motivation}
	
	\newpage
	
	%%%%%%%%%%%%%%%%%%%%%%%%%%%%%%%%%%%%%%%%%%%%%%%%%%%%%
	\subsection*{Séance 1 -- GPIO}
	
	\begin{sessionbox}
		
		\textbf{Date :} Juillet 2026
		
		\vspace{0.3cm}
		
		\textbf{Objectif :}
		
		Faire clignoter la LED intégrée de la carte STM32 Nucleo-F446RE afin de valider la première interaction réelle avec le microcontrôleur.
		
		\vspace{0.3cm}
		
		\textbf{Objectif pédagogique :}
		
		Comprendre le principe d'un GPIO non pas comme une simple fonction logicielle, mais comme un périphérique matériel contrôlé par des registres mémoire. L'objectif était aussi de commencer à comprendre la relation entre le processeur, la carte mémoire, les bus internes, les registres et les broches physiques du microcontrôleur.
		
		\vspace{0.3cm}
		
		\textbf{Réalisation :}
		
		\begin{itemize}
			\item Mise en place du dossier de travail de la séance 1.
			\item Création d'un projet STM32 minimal nommé \texttt{allume\_led}.
			\item Mise en place d'un workflow de développement basé sur VS Code, CMake, Ninja, la toolchain \texttt{arm-none-eabi-gcc} et OpenOCD.
			\item Vérification de la détection du ST-Link intégré de la carte Nucleo-F446RE sous Linux.
			\item Compilation réussie du projet en mode \texttt{Debug}.
			\item Écriture d'un programme de clignotement de LED sans HAL, par accès direct aux registres.
			\item Flash du programme sur la STM32 avec OpenOCD.
			\item Validation expérimentale : la LED LD2 reliée à PA5 clignote correctement.
			\item Mise en place d'un fichier \texttt{.gitignore} propre pour éviter de versionner les fichiers générés.
			\item Commit Git et synchronisation du projet avec GitHub.
		\end{itemize}
		
		\vspace{0.3cm}
		
		\textbf{Choix technique :}
		
		Pour cette séance, le choix a été fait de ne pas utiliser HAL. L'objectif n'était pas seulement d'obtenir un résultat fonctionnel, mais de comprendre le fonctionnement interne du microcontrôleur. Le programme a donc été écrit en manipulant directement les adresses mémoire et les registres nécessaires.
		
		\vspace{0.3cm}
		
		\textbf{Problèmes rencontrés :}
		
		\begin{itemize}
			\item Le bouton \textit{Launch STM32CubeMX} de l'extension VS Code ne trouvait pas STM32CubeMX, car STM32CubeMX n'était pas installé ou pas présent dans le chemin système.
			\item Le projet créé était un projet minimal et non un projet STM32Cube/HAL.
			\item La première configuration CMake échouait car Ninja n'était pas installé.
			\item La compilation échouait ensuite car la toolchain ARM \texttt{arm-none-eabi-gcc} n'était pas disponible dans le \texttt{PATH}.
			\item Certaines fonctions graphiques avancées de l'extension STM32 pour VS Code n'étaient pas directement disponibles car le projet n'était pas reconnu comme un projet STM32Cube complet.
		\end{itemize}
		
		\vspace{0.3cm}
		
		\textbf{Solutions trouvées :}
		
		\begin{itemize}
			\item Utilisation du terminal comme workflow principal pour compiler et flasher le programme.
			\item Installation de Ninja, CMake et de la toolchain ARM.
			\item Compilation avec \texttt{cmake --preset Debug} puis \texttt{cmake --build --preset Debug}.
			\item Flash de la carte avec OpenOCD via l'interface ST-Link.
			\item Décision de conserver ce premier projet comme projet pédagogique minimal, centré sur la compréhension des registres et du fonctionnement matériel.
		\end{itemize}
		
		\vspace{0.3cm}
		
		\textbf{Notions comprises :}
		
		\begin{itemize}
			\item Un GPIO est un périphérique matériel interne au microcontrôleur.
			\item Le processeur ne commande pas directement une LED : il écrit dans des registres mémoire.
			\item La STM32 utilise le principe du \textit{Memory-Mapped I/O} : les périphériques sont accessibles à travers des adresses mémoire.
			\item \texttt{GPIOA\_BASE} représente l'adresse de base du périphérique GPIOA.
			\item Les registres d'un périphérique sont situés à des offsets fixes à partir de son adresse de base.
			\item Le registre \texttt{MODER} configure le mode des broches GPIO.
			\item Le registre \texttt{ODR} permet de modifier l'état logique des sorties.
			\item Avant d'utiliser GPIOA, il faut activer son horloge via le registre \texttt{RCC\_AHB1ENR}.
			\item AHB signifie \textit{Advanced High-performance Bus}. C'est un bus interne rapide reliant le processeur à certains périphériques.
			\item Le mot-clé \texttt{volatile} est essentiel pour empêcher le compilateur d'optimiser les accès aux registres matériels.
			\item Les opérations binaires comme le décalage, le masquage et le XOR permettent de modifier précisément certains bits d'un registre.
		\end{itemize}
		
		\vspace{0.3cm}
		
		\textbf{Explication synthétique du fonctionnement :}
		
		Le processeur Cortex-M4 écrit dans le registre d'activation d'horloge du RCC afin d'activer le périphérique GPIOA. Ensuite, il configure la broche PA5 en sortie en modifiant les deux bits correspondants dans le registre \texttt{MODER}. Une fois PA5 configurée en sortie, le programme inverse régulièrement le bit 5 du registre \texttt{ODR}. Cette modification logique se traduit matériellement par un changement de tension sur la broche PA5, ce qui provoque l'allumage puis l'extinction de la LED LD2.
		
		\vspace{0.3cm}
		
		\textbf{Méthode d'apprentissage retenue :}
		
		La méthode retenue pour la suite du projet est de comprendre chaque périphérique avant de l'utiliser. Pour chaque nouvelle brique, il faudra identifier le rôle du périphérique, sa position dans la carte mémoire, ses registres principaux, puis écrire le code correspondant. Le but n'est pas seulement d'utiliser une bibliothèque, mais d'apprendre à lire la documentation constructeur et à raisonner directement sur l'architecture du microcontrôleur.
		
		\vspace{0.3cm}
		
		\textbf{Résultat :}
		
		La LED LD2 de la carte Nucleo-F446RE clignote correctement. La chaîne complète de développement est validée : édition du code, compilation, génération de l'exécutable, flash avec ST-Link et exécution sur la carte.
		
		\vspace{0.3cm}
		
		\textbf{Apprentissage principal :}
		
		Le processeur ne commande jamais directement une LED. Il écrit dans des registres mémoire, et ce sont les périphériques matériels qui transforment ces écritures en actions physiques.
		
		\vspace{0.3cm}
		
		\textbf{Prochaine étape :}
		
		Commencer la séance 2 avec l'UART afin de mettre en place un premier outil de debug série. Cette étape permettra d'afficher des messages depuis la STM32 vers le PC et de mieux observer l'état interne du programme pendant les essais.
		
	\end{sessionbox}
	
	
	%%%%%%%%%%%%%%%%%%%%%%%%%%%%%%%%%%%%%%%%%%%%%%%%%%%%%
	\subsection*{Séance 2 -- UART}
	
	\begin{sessionbox}
		
		\textbf{Date :} Juillet 2026
		
		\vspace{0.3cm}
		
		\textbf{Objectif :}
		
		Mettre en place une liaison série entre la STM32 Nucleo-F446RE et le PC afin de disposer d'une console de debug simple, fiable et réutilisable pendant tout le projet.
		
		\vspace{0.3cm}
		
		\textbf{Objectif pédagogique :}
		
		Comprendre comment un périphérique UART est configuré en accès direct aux registres : activation des horloges, configuration d'une broche en fonction alternative, choix du débit, activation de l'émetteur et envoi des données caractère par caractère.
		
		\vspace{0.3cm}
		
		\textbf{Configuration retenue :}
		
		\begin{itemize}
			\item périphérique : \texttt{USART2} ;
			\item broche d'émission : \texttt{PA2} ;
			\item fonction alternative : \texttt{AF7} ;
			\item débit : 115200 bauds ;
			\item format : 8 bits de données, aucune parité, 1 bit de stop, soit \texttt{115200 8N1} ;
			\item terminal Linux : \texttt{picocom} sur \texttt{/dev/ttyACM0}.
		\end{itemize}
		
		\vspace{0.3cm}
		
		\textbf{Réalisation :}
		
		\begin{itemize}
			\item Activation de l'horloge de \texttt{GPIOA} dans \texttt{RCC\_AHB1ENR}.
			\item Activation de l'horloge de \texttt{USART2} dans \texttt{RCC\_APB1ENR}.
			\item Configuration de \texttt{PA2} en mode fonction alternative dans \texttt{GPIOA\_MODER}.
			\item Sélection de la fonction \texttt{AF7} dans \texttt{GPIOA\_AFRL} afin de relier physiquement \texttt{PA2} à \texttt{USART2\_TX}.
			\item Configuration du registre \texttt{USART2\_BRR} avec la valeur \texttt{0x8B} pour obtenir environ 115200 bauds à partir d'une horloge de 16 MHz.
			\item Activation de l'émetteur avec le bit \texttt{TE} puis activation du périphérique avec le bit \texttt{UE} dans \texttt{USART2\_CR1}.
			\item Écriture de la fonction \texttt{uart\_send\_char()} qui attend le bit \texttt{TXE} avant d'écrire dans le registre \texttt{USART2\_DR}.
			\item Écriture de la fonction \texttt{uart\_send\_string()} pour transmettre une chaîne caractère par caractère jusqu'au caractère terminal \texttt{'\textbackslash 0'}.
			\item Écriture de la fonction \texttt{uart\_send\_uint()} pour convertir un entier non signé en caractères décimaux sans utiliser \texttt{printf()}.
			\item Utilisation d'une temporisation simple basée sur une boucle et l'instruction \texttt{nop} pour espacer les messages.
			\item Compilation avec CMake, flash avec OpenOCD et observation des messages dans \texttt{picocom}.
		\end{itemize}
		
		\vspace{0.3cm}
		
		\textbf{Résultats obtenus :}
		
		Le premier test a permis d'afficher correctement :
		
		\begin{center}
			\texttt{Hello Robot}
		\end{center}
		
		Le programme a ensuite été enrichi afin d'afficher une suite de messages numérotés :
		
		\begin{center}
			\texttt{Message numero 1}\\
			\texttt{Message numero 2}\\
			\texttt{Message numero 3}
		\end{center}
		
		La chaîne complète de communication a ainsi été validée :
		
		\begin{center}
			STM32 $\rightarrow$ USART2 $\rightarrow$ PA2 $\rightarrow$ ST-Link $\rightarrow$ USB $\rightarrow$ \texttt{/dev/ttyACM0} $\rightarrow$ \texttt{picocom}
		\end{center}
		
		\vspace{0.3cm}
		
		\textbf{Organisation logicielle :}
		
		Le code UART a été séparé du programme principal afin de commencer à construire une architecture modulaire :
		
		\begin{center}
			\begin{tabular}{l}
				\texttt{Core/Inc/uart.h} : déclarations de l'interface publique ;\\
				\texttt{Core/Src/uart.c} : configuration des registres et fonctions d'émission ;\\
				\texttt{Core/Src/main.c} : logique principale de l'application.
			\end{tabular}
		\end{center}
		
		Le fichier \texttt{main.c} ne manipule donc plus directement les registres de l'USART. Il utilise uniquement les fonctions \texttt{uart\_init()}, \texttt{uart\_send\_char()}, \texttt{uart\_send\_string()} et \texttt{uart\_send\_uint()}.
		
		\vspace{0.3cm}
		
		\textbf{Choix concernant \texttt{printf()} :}
		
		La redirection de \texttt{printf()} vers l'UART a volontairement été reportée. L'objectif de cette séance était de comprendre les mécanismes fondamentaux de transmission et la conversion manuelle des nombres, sans masquer le fonctionnement derrière la bibliothèque standard.
		
		\vspace{0.3cm}
		
		\textbf{Problèmes rencontrés :}
		
		\begin{itemize}
			\item Le terminal série \texttt{picocom} n'était pas initialement installé sur le système.
			\item Le message \texttt{Hello Robot} n'est envoyé qu'une fois au démarrage : si le terminal est ouvert après l'exécution du programme, le message n'est pas visible.
		\end{itemize}
		
		\vspace{0.3cm}
		
		\textbf{Solutions trouvées :}
		
		\begin{itemize}
			\item Installation de \texttt{picocom} puis ouverture du port avec la commande \texttt{picocom -b 115200 /dev/ttyACM0}.
			\item Appui sur le bouton \texttt{RESET} de la carte après l'ouverture du terminal afin de relancer le programme et de recevoir le message initial.
		\end{itemize}
		
		\vspace{0.3cm}
		
		\textbf{Notions comprises :}
		
		\begin{itemize}
			\item \texttt{TX} correspond à la transmission et \texttt{RX} à la réception.
			\item Une communication UART est asynchrone : les deux équipements doivent utiliser le même débit et le même format de trame.
			\item Une broche en fonction alternative est pilotée par un périphérique matériel et non directement par le registre \texttt{ODR} du GPIO.
			\item Le débit est défini par le registre \texttt{BRR} à partir de la fréquence d'horloge du périphérique.
			\item Le bit \texttt{TXE} indique que le registre de transmission peut recevoir un nouveau caractère.
			\item Le registre \texttt{DR} contient la donnée à transmettre.
			\item Une chaîne de caractères est un tableau terminé par \texttt{'\textbackslash 0'}.
			\item Un entier doit être converti en caractères avant de pouvoir être affiché dans un terminal.
			\item La séparation entre interface \texttt{.h}, implémentation \texttt{.c} et programme principal rend le code plus lisible et réutilisable.
		\end{itemize}
		
		\vspace{0.3cm}
		
		\textbf{Apprentissage principal :}
		
		L'UART transforme les données écrites par le processeur dans des registres en une suite de bits transmise sur une broche physique. Grâce à cette liaison, le programme embarqué peut rendre visible son état interne sur le PC.
		
		\vspace{0.3cm}
		
		\textbf{Prochaine étape :}
		
		Utiliser cette console UART pendant la séance 3 pour afficher le résultat de la première communication I2C avec le MPU6050, notamment la valeur de son registre d'identification \texttt{WHO\_AM\_I}.
		
	\end{sessionbox}
	
	%%%%%%%%%%%%%%%%%%%%%%%%%%%%%%%%%%%%%%%%%%%%%%%%%%%%%
	\subsection*{Séance 3 -- IMU : première communication}
	
	\begin{sessionbox}
		
		\textbf{Objectif :} Vérifier que le module MPU6050 communique correctement avec la STM32.
		
		\vspace{0.3cm}
		
		\textbf{Technologie :} I2C.
		
		\vspace{0.3cm}
		
		\textbf{Réalisation :}
		
		Lire le registre d'identification de l'IMU et afficher la réponse via UART.
		
		\vspace{0.3cm}
		
		\textbf{Résultat attendu :}
		
		Confirmer que le capteur répond à l'adresse I2C prévue.
		
		\vspace{0.3cm}
		
		\textbf{Apprentissage :}
		
		Avant d'exploiter un capteur, il faut d'abord valider sa communication électrique et logicielle.
		
	\end{sessionbox}
	
	%%%%%%%%%%%%%%%%%%%%%%%%%%%%%%%%%%%%%%%%%%%%%%%%%%%%%
	\subsection*{Séance 4 -- PWM moteur}
	
	\begin{sessionbox}
		
		\textbf{Objectif :} Générer un signal PWM pour contrôler la vitesse moteur.
		
		\vspace{0.3cm}
		
		\textbf{Technologie :} Timers + PWM.
		
		\vspace{0.3cm}
		
		\textbf{Réalisation :}
		
		Configurer un timer STM32 en mode PWM et faire varier le rapport cyclique.
		
		\vspace{0.3cm}
		
		\textbf{Apprentissage :}
		
		Le PWM permet de contrôler l'énergie moyenne envoyée au moteur.
		
	\end{sessionbox}
	
	%%%%%%%%%%%%%%%%%%%%%%%%%%%%%%%%%%%%%%%%%%%%%%%%%%%%%
	\subsection*{Séance 5 -- Encodeurs}
	
	\begin{sessionbox}
		
		\textbf{Objectif :} Mesurer la vitesse réelle des moteurs.
		
		\vspace{0.3cm}
		
		\textbf{Technologie :} Timer en mode encodeur.
		
		\vspace{0.3cm}
		
		\textbf{Réalisation :}
		
		Connecter les voies A et B d'un encodeur à deux entrées timer et lire la position moteur.
		
		\vspace{0.3cm}
		
		\textbf{Apprentissage :}
		
		L'encodeur permet de passer d'une commande aveugle à une mesure réelle du mouvement.
		
	\end{sessionbox}
	
	%%%%%%%%%%%%%%%%%%%%%%%%%%%%%%%%%%%%%%%%%%%%%%%%%%%%%
	\subsection*{Séance 6 -- Lecture moteur + sens de rotation}
	
	\begin{sessionbox}
		
		\textbf{Objectif :} Contrôler la direction et le sens des moteurs.
		
		\vspace{0.3cm}
		
		\textbf{Technologie :} GPIO + PWM + Driver moteur TB6612FNG.
		
		\vspace{0.3cm}
		
		\textbf{Réalisation :}
		
		\begin{itemize}
			\item inversion du sens de rotation ;
			\item test avant/arrière ;
			\item validation du pont en H TB6612FNG.
		\end{itemize}
		
		\vspace{0.3cm}
		
		\textbf{Problème possible :}
		
		Le moteur tourne dans le mauvais sens.
		
		\vspace{0.3cm}
		
		\textbf{Solution :}
		
		Inversion des broches de direction IN1/IN2 ou correction logicielle du signe de commande.
		
		\vspace{0.3cm}
		
		\textbf{Apprentissage :}
		
		Un moteur ne se contrôle pas seulement en vitesse, mais aussi en direction.
		
	\end{sessionbox}
	
	%%%%%%%%%%%%%%%%%%%%%%%%%%%%%%%%%%%%%%%%%%%%%%%%%%%%%
	\subsection*{Séance 7 -- Lecture IMU brute}
	
	\begin{sessionbox}
		
		\textbf{Objectif :} Lire les données brutes de l'IMU.
		
		\vspace{0.3cm}
		
		\textbf{Technologie :} I2C + MPU6050.
		
		\vspace{0.3cm}
		
		\textbf{Réalisation :}
		
		\begin{itemize}
			\item lecture des registres d'accélération ;
			\item lecture du gyroscope ;
			\item affichage UART.
		\end{itemize}
		
		\vspace{0.3cm}
		
		\textbf{Problème attendu :}
		
		Valeurs instables et bruitées.
		
		\vspace{0.3cm}
		
		\textbf{Apprentissage :}
		
		Les capteurs bruts ne sont jamais directement exploitables pour le contrôle.
		
	\end{sessionbox}
	
	%%%%%%%%%%%%%%%%%%%%%%%%%%%%%%%%%%%%%%%%%%%%%%%%%%%%%
	\subsection*{Séance 8 -- Filtre complémentaire}
	
	\begin{sessionbox}
		
		\textbf{Objectif :} Obtenir un angle stable.
		
		\vspace{0.3cm}
		
		\textbf{Technologie :} Fusion accéléromètre + gyroscope.
		
		\vspace{0.3cm}
		
		\textbf{Réalisation :}
		
		\begin{itemize}
			\item implémentation du filtre complémentaire ;
			\item estimation de l'angle de tangage ;
			\item stabilisation des mesures.
		\end{itemize}
		
		\vspace{0.3cm}
		
		\textbf{Problème attendu :}
		
		Dérive gyroscopique et bruit de l'accéléromètre.
		
		\vspace{0.3cm}
		
		\textbf{Solution :}
		
		Pondérer l'information du gyroscope et de l'accéléromètre.
		
		\vspace{0.3cm}
		
		\textbf{Apprentissage :}
		
		La fusion de capteurs est une base des systèmes autonomes.
		
	\end{sessionbox}
	
	%%%%%%%%%%%%%%%%%%%%%%%%%%%%%%%%%%%%%%%%%%%%%%%%%%%%%
	\subsection*{Séance 9 -- Test moteur en boucle ouverte}
	
	\begin{sessionbox}
		
		\textbf{Objectif :} Tester les moteurs sans contrôle fermé.
		
		\vspace{0.3cm}
		
		\textbf{Technologie :} PWM fixe.
		
		\vspace{0.3cm}
		
		\textbf{Réalisation :}
		
		\begin{itemize}
			\item test à vitesse constante ;
			\item test sous différentes charges ;
			\item observation du comportement mécanique.
		\end{itemize}
		
		\vspace{0.3cm}
		
		\textbf{Problème attendu :}
		
		La vitesse varie selon la charge et l'état de la batterie.
		
		\vspace{0.3cm}
		
		\textbf{Apprentissage :}
		
		La boucle ouverte est insuffisante pour un système dynamique précis.
		
	\end{sessionbox}
	
	%%%%%%%%%%%%%%%%%%%%%%%%%%%%%%%%%%%%%%%%%%%%%%%%%%%%%
	\subsection*{Séance 10 -- Boucle de contrôle vitesse}
	
	\begin{sessionbox}
		
		\textbf{Objectif :} Stabiliser la vitesse moteur.
		
		\vspace{0.3cm}
		
		\textbf{Technologie :} PID vitesse + encodeurs.
		
		\vspace{0.3cm}
		
		\textbf{Réalisation :}
		
		\begin{itemize}
			\item implémentation d'un PID de vitesse ;
			\item lecture encodeur ;
			\item ajustement automatique du PWM.
		\end{itemize}
		
		\vspace{0.3cm}
		
		\textbf{Problème attendu :}
		
		Oscillations ou réponse trop lente.
		
		\vspace{0.3cm}
		
		\textbf{Solution :}
		
		Régler progressivement les gains, en commençant par le gain proportionnel.
		
		\vspace{0.3cm}
		
		\textbf{Apprentissage :}
		
		Un PID mal réglé crée de l'instabilité au lieu de la corriger.
		
	\end{sessionbox}
	
	%%%%%%%%%%%%%%%%%%%%%%%%%%%%%%%%%%%%%%%%%%%%%%%%%%%%%
	\subsection*{Séance 11 -- Boucle d'équilibre en simulation}
	
	\begin{sessionbox}
		
		\textbf{Objectif :} Tester le contrôle d'équilibre avant l'intégration physique.
		
		\vspace{0.3cm}
		
		\textbf{Technologie :} Simulation logicielle d'un pendule inversé simplifié.
		
		\vspace{0.3cm}
		
		\textbf{Réalisation :}
		
		\begin{itemize}
			\item modélisation simplifiée de l'angle ;
			\item test du PID angle ;
			\item validation du comportement attendu.
		\end{itemize}
		
		\vspace{0.3cm}
		
		\textbf{Apprentissage :}
		
		Toujours tester un contrôle avant de l'appliquer directement à un système instable réel.
		
	\end{sessionbox}
	
	%%%%%%%%%%%%%%%%%%%%%%%%%%%%%%%%%%%%%%%%%%%%%%%%%%%%%
	\subsection*{Séance 12 -- Construction du châssis}
	
	\begin{sessionbox}
		
		\textbf{Objectif :} Construire la base mécanique du robot.
		
		\vspace{0.3cm}
		
		\textbf{Technologie :} Mécanique + intégration électronique.
		
		\vspace{0.3cm}
		
		\textbf{Réalisation :}
		
		\begin{itemize}
			\item assemblage de la structure du robot ;
			\item fixation des moteurs ;
			\item positionnement de la batterie pour maîtriser le centre de gravité ;
			\item installation de la STM32 et du driver moteur ;
			\item câblage initial.
		\end{itemize}
		
		\vspace{0.3cm}
		
		\textbf{Points critiques :}
		
		\begin{itemize}
			\item équilibre mécanique ;
			\item rigidité du châssis ;
			\item alignement des roues ;
			\item fixation stable de l'IMU.
		\end{itemize}
		
		\vspace{0.3cm}
		
		\textbf{Apprentissage :}
		
		La stabilité d'un robot commence en mécanique, pas en code.
		
	\end{sessionbox}
	
	%%%%%%%%%%%%%%%%%%%%%%%%%%%%%%%%%%%%%%%%%%%%%%%%%%%%%
	\subsection*{Séance 13 -- Première tentative d'équilibre}
	
	\begin{sessionbox}
		
		\textbf{Objectif :} Faire tenir le robot debout.
		
		\vspace{0.3cm}
		
		\textbf{Technologie :} IMU + PID angle + PWM.
		
		\vspace{0.3cm}
		
		\textbf{Réalisation :}
		
		\begin{itemize}
			\item boucle de contrôle à 500--1000 Hz ;
			\item lecture IMU en temps réel ;
			\item correction moteur selon l'angle.
		\end{itemize}
		
		\vspace{0.3cm}
		
		\textbf{Résultat probable :}
		
		Oscillations rapides et instabilité initiale.
		
		\vspace{0.3cm}
		
		\textbf{Apprentissage :}
		
		Un système instable doit être réglé, pas abandonné.
		
	\end{sessionbox}
	
	%%%%%%%%%%%%%%%%%%%%%%%%%%%%%%%%%%%%%%%%%%%%%%%%%%%%%
	\subsection*{Séance 14 -- Stabilisation progressive}
	
	\begin{sessionbox}
		
		\textbf{Objectif :} Améliorer la stabilité du robot.
		
		\vspace{0.3cm}
		
		\textbf{Actions :}
		
		\begin{itemize}
			\item réglage progressif du PID ;
			\item réduction du bruit IMU ;
			\item amélioration de la fréquence de boucle ;
			\item observation du comportement mécanique.
		\end{itemize}
		
		\vspace{0.3cm}
		
		\textbf{Résultat attendu :}
		
		Réduction des oscillations et maintien de l'équilibre pendant plusieurs secondes.
		
		\vspace{0.3cm}
		
		\textbf{Apprentissage :}
		
		Le contrôle est un processus itératif.
		
	\end{sessionbox}
	
	%%%%%%%%%%%%%%%%%%%%%%%%%%%%%%%%%%%%%%%%%%%%%%%%%%%%%
	\subsection*{Séance 15 -- Déplacement contrôlé}
	
	\begin{sessionbox}
		
		\textbf{Objectif :} Faire avancer et reculer le robot sans tomber.
		
		\vspace{0.3cm}
		
		\textbf{Technologie :} Contrôle angle + consigne vitesse.
		
		\vspace{0.3cm}
		
		\textbf{Réalisation :}
		
		\begin{itemize}
			\item ajout d'une consigne de vitesse ;
			\item couplage entre équilibre et déplacement ;
			\item test de stabilité dynamique.
		\end{itemize}
		
		\vspace{0.3cm}
		
		\textbf{Apprentissage :}
		
		Stabilité et mouvement forment un système couplé.
		
	\end{sessionbox}
	
	%%%%%%%%%%%%%%%%%%%%%%%%%%%%%%%%%%%%%%%%%%%%%%%%%%%%%
	\subsection*{Séance 16 -- Préparation autonomie}
	
	\begin{sessionbox}
		
		\textbf{Objectif :} Préparer une navigation simple d'un point A vers un point B.
		
		\vspace{0.3cm}
		
		\textbf{Technologie :} Odométrie + contrôle position.
		
		\vspace{0.3cm}
		
		\textbf{Réalisation :}
		
		\begin{itemize}
			\item estimation de la distance parcourue ;
			\item correction de trajectoire ;
			\item test d'un mouvement simple sur une distance donnée.
		\end{itemize}
		
		\vspace{0.3cm}
		
		\textbf{Apprentissage :}
		
		Un robot stable est une base pour l'autonomie.
		
	\end{sessionbox}
	
	% ================================
	% CHAPITRE 5
	% ================================
	\section{Les composants}
	
	\subsection*{STM32 Nucleo-F446RE}
	
	\begin{componentbox}
		\begin{itemize}
			\item Microcontrôleur STM32F446RE.
			\item Coeur ARM Cortex-M4.
			\item Timers adaptés au PWM et aux encodeurs.
			\item Interfaces UART, SPI, I2C.
			\item ADC, DMA et interruptions.
			\item ST-Link intégré pour programmation et debug.
		\end{itemize}
	\end{componentbox}
	
	\subsection*{IMU MPU6050}
	
	\begin{componentbox}
		\begin{itemize}
			\item Accéléromètre 3 axes.
			\item Gyroscope 3 axes.
			\item Communication I2C.
			\item Utilisé pour estimer l'angle du robot.
		\end{itemize}
	\end{componentbox}
	
	\subsection*{Moteurs avec encodeurs}
	
	\begin{componentbox}
		\begin{itemize}
			\item Moteurs DC avec réduction mécanique.
			\item Encodeurs quadrature voies A et B.
			\item Mesure de vitesse et de position.
			\item Utilisés pour le contrôle de vitesse et l'odométrie.
		\end{itemize}
	\end{componentbox}
	
	\subsection*{Driver moteur TB6612FNG}
	
	\begin{componentbox}
		\begin{itemize}
			\item Double pont en H.
			\item Contrôle de deux moteurs DC.
			\item Commande par GPIO et PWM.
			\item Alimentation moteur séparée de la logique.
		\end{itemize}
	\end{componentbox}
	
	\subsection*{Alimentation}
	
	\begin{componentbox}
		\begin{itemize}
			\item Batterie LiPo 2S 7,4 V.
			\item Convertisseur Buck LM2596 pour obtenir une tension logique stable.
			\item Interrupteur ON/OFF sur la ligne positive.
			\item Connecteurs Deans pour la puissance.
		\end{itemize}
	\end{componentbox}
	
	% ================================
	% CHAPITRE 6
	% ================================
	\section{Notions théoriques}
	
	\subsection*{Memory-Mapped I/O : carte mémoire, périphériques et registres}
	
	Dans un microcontrôleur STM32, l'espace d'adressage ne contient pas uniquement la mémoire dans laquelle sont stockées les variables. Une partie des adresses est réservée aux périphériques matériels : GPIO, UART, I2C, timers, ADC ou encore RCC. Ce principe est appelé \textit{Memory-Mapped I/O}.
	
	Pour le processeur, lire ou écrire un registre matériel ressemble donc à lire ou écrire une case mémoire. La différence est que cette opération produit un effet physique ou permet de lire l'état d'un périphérique.
	
	\begin{center}
		\begin{tabular}{>{\bfseries}p{4cm}p{9cm}}
			\toprule
			Élément & Signification \\
			\midrule
			Carte mémoire & Organisation de tout l'espace d'adresses vu par le processeur. \\
			Périphérique & Bloc matériel, par exemple GPIOA ou USART2. \\
			Adresse de base & Première adresse réservée au périphérique. \\
			Offset & Position d'un registre par rapport à l'adresse de base. \\
			Registre & Petite zone mémoire, souvent de 32 bits, qui configure ou décrit le périphérique. \\
			Bit ou champ de bits & Partie précise d'un registre associée à une fonction. \\
			\bottomrule
		\end{tabular}
	\end{center}
	
	La règle essentielle est :
	
	\begin{center}
		\textbf{Adresse du registre = adresse de base du périphérique + offset du registre}
	\end{center}
	
	Exemple avec GPIOA :
	
	\begin{itemize}
		\item \texttt{GPIOA\_BASE = 0x40020000} ;
		\item l'offset de \texttt{MODER} est \texttt{0x00} ;
		\item l'adresse de \texttt{GPIOA\_MODER} est donc \texttt{0x40020000} ;
		\item l'offset de \texttt{ODR} est \texttt{0x14} ;
		\item l'adresse de \texttt{GPIOA\_ODR} est donc \texttt{0x40020014}.
	\end{itemize}
	
	Exemple avec USART2 :
	
	\begin{center}
		\begin{tabular}{>{\bfseries}p{3.5cm}p{3cm}p{5cm}}
			\toprule
			Registre & Offset & Rôle principal \\
			\midrule
			\texttt{USART2\_SR} & \texttt{0x00} & Indique l'état du périphérique, par exemple \texttt{TXE}. \\
			\texttt{USART2\_DR} & \texttt{0x04} & Contient la donnée à transmettre ou reçue. \\
			\texttt{USART2\_BRR} & \texttt{0x08} & Configure le débit en bauds. \\
			\texttt{USART2\_CR1} & \texttt{0x0C} & Active et configure les fonctions principales de l'USART. \\
			\bottomrule
		\end{tabular}
	\end{center}
	
	Une définition comme :
	
	\begin{center}
		\texttt{\#define USART2\_DR (*(volatile uint32\_t *)(USART2\_BASE + 0x04UL))}
	\end{center}
	
	peut être comprise en quatre étapes :
	
	\begin{enumerate}
		\item calculer l'adresse du registre avec la base et l'offset ;
		\item considérer cette adresse comme un pointeur vers un entier de 32 bits ;
		\item utiliser \texttt{volatile} pour imposer un véritable accès matériel ;
		\item déréférencer le pointeur avec \texttt{*} pour lire ou modifier le registre.
	\end{enumerate}
	
	\begin{warningbox}
		Une adresse incorrecte ou un mauvais bit peut configurer un autre périphérique ou produire un comportement imprévisible. Les adresses, offsets et positions des bits doivent toujours être vérifiés dans le manuel de référence et la datasheet du microcontrôleur.
	\end{warningbox}
	
	\subsection*{Registres et opérations sur les bits}
	
	Un registre contient plusieurs fonctions indépendantes. Il ne faut donc généralement pas remplacer toute sa valeur lorsqu'on souhaite modifier un seul bit.
	
	Les opérations les plus utilisées sont :
	
	\begin{itemize}
		\item activer un bit : \verb!registre |= (1UL << n)! ;
		\item désactiver un bit : \verb!registre &= ~(1UL << n)! ;
		\item tester un bit : \verb!registre & (1UL << n)! ;
		\item inverser un bit : \verb!registre ^= (1UL << n)! ;
		\item modifier un champ de plusieurs bits : effacer d'abord le champ avec un masque, puis écrire la nouvelle valeur.
	\end{itemize}
	
	Cette méthode permet de modifier uniquement la fonction recherchée sans perturber les autres bits du registre.
	
	\subsection*{RCC et bus internes}
	
	Le RCC, ou \textit{Reset and Clock Control}, distribue les horloges aux différents périphériques. Un périphérique dont l'horloge n'est pas activée ne peut pas fonctionner correctement, même si ses registres sont configurés.
	
	Les périphériques sont reliés au processeur par plusieurs bus internes :
	
	\begin{itemize}
		\item AHB1 pour des périphériques rapides comme les GPIO ;
		\item APB1 pour plusieurs périphériques comme USART2, I2C et certains timers ;
		\item APB2 pour d'autres périphériques rapides comme certains USART et timers.
	\end{itemize}
	
	La démarche habituelle est donc : identifier le bus du périphérique, activer son horloge dans le registre RCC correspondant, puis configurer ses propres registres.
	
	\subsection*{GPIO}
	
	Les GPIO sont des broches numériques configurables. Une broche peut notamment fonctionner comme entrée, sortie, fonction alternative ou entrée analogique. Le registre \texttt{MODER} choisit ce mode.
	
	En sortie GPIO, le programme commande directement le niveau logique grâce au registre \texttt{ODR}. En fonction alternative, la broche est confiée à un autre périphérique matériel, par exemple un USART ou un timer.
	
	Les GPIO seront utilisés pour les LEDs, les boutons, les signaux de direction des moteurs et plusieurs fonctions alternatives.
	
	\subsection*{UART}
	
	L'UART est une communication série asynchrone. Elle utilise principalement une ligne \texttt{TX} pour transmettre et une ligne \texttt{RX} pour recevoir. Comme aucune horloge n'est transmise sur un fil séparé, les deux équipements doivent partager le même débit et le même format de trame.
	
	La configuration \texttt{115200 8N1} signifie : 115200 symboles par seconde, 8 bits de données, aucune parité et 1 bit de stop.
	
	Dans ce projet, USART2 transmet sur PA2. Les registres principaux utilisés sont :
	
	\begin{itemize}
		\item \texttt{BRR} pour le débit ;
		\item \texttt{CR1} pour activer l'émetteur et le périphérique ;
		\item \texttt{SR} pour consulter l'état de la transmission ;
		\item \texttt{DR} pour écrire le caractère à transmettre.
	\end{itemize}
	
	L'UART sert au debug et à la télémétrie. Il permet d'observer les mesures, les états et les erreurs internes du robot depuis un terminal PC.
	
	\subsection*{I2C}
	
	L'I2C est un bus série synchrone utilisant deux lignes : \texttt{SDA} pour les données et \texttt{SCL} pour l'horloge. Plusieurs périphériques peuvent partager les mêmes lignes, chaque composant étant identifié par une adresse.
	
	Dans ce projet, l'I2C permettra à la STM32 de sélectionner un registre du MPU6050, puis de lire ou d'écrire sa valeur. Des résistances de tirage vers le niveau haut sont nécessaires sur SDA et SCL ; elles sont souvent déjà présentes sur les modules MPU6050.
	
	\subsection*{PWM}
	
	Le PWM, ou modulation de largeur d'impulsion, est un signal numérique périodique dont on fait varier le rapport cyclique. Un rapport cyclique plus élevé augmente la durée pendant laquelle le signal reste à l'état haut et donc la puissance moyenne appliquée au moteur.
	
	Le PWM sera généré par un timer STM32 puis transmis au driver TB6612FNG. Il ne modifie pas directement la tension de la batterie : il découpe rapidement l'alimentation afin de contrôler l'énergie moyenne reçue par le moteur.
	
	\subsection*{Encodeurs}
	
	Les encodeurs produisent des impulsions lorsque le moteur tourne. Avec deux voies A et B déphasées, ils permettent de déterminer le sens de rotation ainsi que le déplacement relatif.
	
	En comptant les impulsions pendant une durée connue, le programme peut estimer la vitesse. En cumulant les impulsions, il peut estimer la position ou la distance parcourue.
	
	\subsection*{Filtre complémentaire}
	
	L'accéléromètre fournit une estimation de l'inclinaison stable à long terme mais sensible aux vibrations et aux accélérations du robot. Le gyroscope mesure rapidement les variations angulaires mais son intégration dérive progressivement.
	
	Le filtre complémentaire combine les deux informations : le gyroscope décrit principalement les variations rapides, tandis que l'accéléromètre corrige lentement la dérive. Le résultat est une estimation d'angle plus stable pour la boucle d'équilibre.
	
	\subsection*{PID}
	
	Le PID corrige l'erreur entre une consigne et une mesure :
	
	\begin{itemize}
		\item le terme proportionnel réagit à l'erreur actuelle ;
		\item le terme intégral corrige une erreur persistante ;
		\item le terme dérivé anticipe les variations rapides de l'erreur.
	\end{itemize}
	
	Dans ce projet, un PID sera utilisé pour contrôler l'angle du robot et un autre pourra être utilisé pour la vitesse des moteurs. Les gains devront être réglés progressivement afin d'éviter les oscillations et l'instabilité.
	
	% ================================
	% CHAPITRE 7
	% ================================
	\section{Architecture du robot}
	
	\subsection*{Architecture matérielle}
	
	\begin{center}
		\begin{tabular}{>{\bfseries}p{4cm}p{9cm}}
			\toprule
			Élément & Rôle \\
			\midrule
			STM32 Nucleo-F446RE & Contrôle temps réel, lecture capteurs, PID, PWM \\
			MPU6050 & Mesure d'accélération et vitesse angulaire \\
			TB6612FNG & Interface de puissance entre STM32 et moteurs \\
			Moteurs + encodeurs & Action mécanique + mesure du mouvement \\
			LiPo 2S & Source d'énergie principale \\
			Buck LM2596 & Conversion vers une tension logique stable \\
			Châssis PVC expansé & Structure mécanique du robot \\
			\bottomrule
		\end{tabular}
	\end{center}
	
	\subsection*{Principe de fonctionnement}
	
	\begin{itemize}
		\item L'IMU mesure l'inclinaison du robot.
		\item La STM32 estime l'angle avec un filtre complémentaire.
		\item Le PID calcule la correction nécessaire.
		\item La STM32 génère les signaux PWM.
		\item Le driver TB6612FNG applique la puissance aux moteurs.
		\item Les encodeurs mesurent le mouvement réel.
	\end{itemize}
	
	% ================================
	% CHAPITRE 8
	% ================================
	\section{Les erreurs et leçons apprises}
	
	\subsection*{Erreur 1 -- PID instable}
	
	\textbf{Cause possible :} gain proportionnel trop élevé.
	
	\textbf{Solution :} réduire le gain proportionnel puis ajouter progressivement les autres termes.
	
	\textbf{Leçon :} régler P avant I et D.
	
	\subsection*{Erreur 2 -- Robot qui tremble}
	
	\textbf{Cause possible :} bruit IMU, câbles moteurs trop proches des signaux, châssis flexible.
	
	\textbf{Solution :} filtrage, séparation des câbles puissance/signaux, fixation rigide de l'IMU.
	
	\textbf{Leçon :} un problème de stabilité n'est pas toujours un problème de code.
	
	\subsection*{Erreur 3 -- Reset de la STM32}
	
	\textbf{Cause possible :} chute de tension lors de l'appel de courant des moteurs.
	
	\textbf{Solution :} alimentation logique séparée via Buck Converter et masse commune propre.
	
	\textbf{Leçon :} l'alimentation est une partie critique du système.
	
	% ================================
	% CHAPITRE 9
	% ================================
	\section{Améliorations futures}
	
	\begin{itemize}
		\item ESP32 pour Wi-Fi et télémétrie.
		\item Interface web pour visualiser les données en temps réel.
		\item Réglage PID à distance.
		\item Écran OLED pour afficher l'état du robot.
		\item Capteurs de distance VL53L0X.
		\item Caméra.
		\item Navigation autonome.
		\item SLAM simple.
		\item Passage à un châssis imprimé en 3D ou aluminium.
	\end{itemize}
	
	% ================================
	% CHAPITRE 10
	% ================================
	\section{Matériel et plan d'achat}
	
	\subsection*{Matériel acheté pour la version 1}
	
	\begin{longtable}{p{6cm}p{7.5cm}}
		\toprule
		\textbf{Composant} & \textbf{Rôle} \\
		\midrule
		STM32 Nucleo-F446RE & Carte principale temps réel \\
		MPU6050 & Mesure inertielle \\
		TB6612FNG & Driver moteur \\
		Deux motoréducteurs avec encodeurs & Actionnement et mesure de mouvement \\
		LiPo 2S 7,4 V -- 2200 mAh & Batterie principale \\
		Chargeur LiPo avec équilibrage & Recharge sécurisée \\
		Buck LM2596 & Conversion tension logique \\
		Connecteurs Deans + câbles silicone & Distribution puissance \\
		Interrupteur ON/OFF & Coupure alimentation \\
		Multimètre & Diagnostic électrique \\
		Cutter & Fabrication châssis \\
		PVC expansé & Structure mécanique \\
		Fils Dupont et câbles & Connexions logiques et puissance \\
		\bottomrule
	\end{longtable}
	
	\subsection*{Budget}
	
	Le coût total de la version initiale du projet est estimé à :
	
	\begin{center}
		\Huge\textbf{200~\texteuro{}}
	\end{center}
	
	\subsection*{Achats reportés}
	
	\begin{itemize}
		\item ESP32 pour Wi-Fi ;
		\item écran OLED ;
		\item capteurs de distance ;
		\item caméra ;
		\item châssis avancé.
	\end{itemize}
	
	% ================================
	% CONCLUSION
	% ================================
	\section*{Conclusion}
	\addcontentsline{toc}{section}{Conclusion}
	
	\begin{center}
		\Huge
		\textbf{Pourvu qu'aujourd'hui soit mieux qu'hier}
		
		\vspace{0.8cm}
		
		\Large
		Je ne cherche pas à être parfait.
		
		Je cherche simplement à progresser.
		
		Chaque ligne de code écrite.
		
		Chaque bug corrigé.
		
		Chaque composant compris.
		
		Chaque test réalisé.
		
		Me rapproche un peu plus de l'ingénieur que je veux devenir.
	\end{center}
	
	\vfill
	
	\begin{center}
		\Large
		\textbf{Le robot n'est pas la destination.}
		
		\vspace{0.3cm}
		
		\textbf{Le véritable projet, c'est moi.}
	\end{center}
	
\end{document}